\subsection{Домашние задачи на шестую неделю}
\begin{problem}
    Пусть $p \in (0, 1)$ и $\xi$~---~дискретная случайная величина с $\P(\xi = k) = -\dfrac{1}{\ln 1 - p}\dfrac{p^k}{k}$.
    Найдите $\Ex\xi$ и $\D\xi$ с помощью производящей функции.
\end{problem}
\begin{solution}
    Найдём производящую функцию распределения:
    \[
        G(z) = \sum\limits_{k = 1}^{+\infty} -\dfrac{1}{\ln 1 - p} \dfrac{(pz)^k}{k} = -\dfrac{1}{\ln 1 - p}\sum\limits_{k = 1}^{+\infty} \dfrac{(pz)^k}{k} = -\dfrac{-\ln 1 - pz}{\ln 1 - p} = \dfrac{\ln 1 - pz}{\ln 1 - p}.
    \]
    Теперь выразим $\Ex \xi$ и $\D \xi$ через производящую функцию:
    \[
        \Ex \xi = G'(1) =  \dfrac{-p}{1 - p}\cdot \dfrac{1}{\ln 1 - p} = \dfrac{-p}{(1 - p)\ln 1 - p}.
    \]
    И
    \[
        \D \xi = G''(1) + G'(1) - (G'(1))^2 = \dfrac{-p^2}{(1 - p)^2\ln 1 - p} - \dfrac{p}{(1 - p)\ln 1 - p} - \dfrac{p^2}{(1 - p)^2\ln^2 1- p}.
    \]
\end{solution}
\begin{problem}
    Пусть $\xi$~---~целочисленная случайная величина, $\phi_\xi(t)$~---~её характеристическая функция. Доказать, что
    \[
        \P(\xi = n) = \dfrac{1}{2\pi}\int_{-\pi}^{\pi} e^{-int}\phi_{\xi}(t)dt, n \in \Z.
    \]
\end{problem}
\begin{solution}
    По определению
    \[
        \phi_{\xi}(t) = \Ex e^{it\xi} = \sum\limits_{n \in \Z} e^{int}\P(\xi = n)
    \]
    В силу ортогональности системы комплексных экспонент на $[-\pi, \pi]$ и условия нормировки для $\sum_\Z \P(\xi = n) = 1$ после подстановки выражения характеристической функции в интеграл и перестановки интеграла с суммой (по сути, применение теоремы Фубини) мы получаем что не зануляется только $n$-й член и таким образом получается искомая формула.
\end{solution}
\begin{problem}
    Доказать, что $\Gamma$-распределение $\Gamma(\lambda, k), k \geq 1$ с плотностью \[p(x) = \dfrac{1}{\lambda^k \Gamma(k)}x^{k - 1}e^{-x/\lambda}I_{\R_+}(x)\]
    устойчиво относительно сложения, то есть сумма двух независимых $X \sim \Gamma(\lambda, k_X)$ и $Y \sim \Gamma(\lambda, k_Y)$ распределена как $\Gamma(\lambda, k_X + k_Y)$.
\end{problem}
\begin{solution}
    Найдём характеристическую функцию $\eta \sim \Gamma(\lambda, k)$.
    По определению, $\phi_\eta (t) = \Ex e^{it\eta}$ и тогда
    \begin{multline*}
        \dfrac{d}{dt} \phi_{\eta}(t) = \dfrac{d}{dt}\dfrac{1}{\lambda^k \Gamma(k)}\int_0^{+\infty} x^{k - 1}e^{-x/\lambda}e^{itx}dx = \dfrac{1}{\lambda^k \Gamma(k)} \int_0^{+\infty} ix^k e^{x(it - 1/\lambda)}dx = \\ =
        \dfrac{i}{it - 1/\lambda} \cdot \dfrac{1}{\lambda^k \Gamma(k)}\biggr(x^ke^{x(it - 1/\lambda)}\biggr|_0^{+\infty} - k\phi_\eta(t)\biggr) = -\dfrac{k\phi_\eta(t)}{(t + i/\lambda)\lambda^k \Gamma(k)}
    \end{multline*}
    Тогда $\phi_\eta(t)$ получается как решение дифференциального уравнения и имеет вид
    \[
        \phi_{\eta}(t) = c(1/\lambda - it)^{-k}.
    \]
    И из того, что $\phi_\eta(0) = 1$ следует, что $c = 1/\lambda^k$.
    Из вида характеристической функции и независимости случайных величин очевидным образом следует утверждение задачи.
\end{solution}
\begin{problem}
    Вывести характеристические функции для распределения Бернулли, биномиального и распределения Пуассона.
\end{problem}
\begin{solution} По определению $\phi_\xi(t) = \Ex e^{it\xi}$. Найдём эти производящие функции:
    \begin{itemize}
        \item Пусть $X \sim Be(p)$ и тогда
        \[
            \phi_X(t) = \Ex e^{itX} = pe^{it \cdot 1} + (1 - p)e^{it \cdot 0} = (1 - p) + pe^{it}.
        \]
        \item Пусть $X \sim Binom(n, p)$ и тогда
        \[
            \phi_X(t) = \Ex e^{itX} = \sum\limits_{k = 0}^n \binom{n}{k}e^{ikt}p^k(1 - p)^{n - k} = (pe^{it} + (1 - p))^n.
        \]
        \item Пусть $X \sim Poiss(\lambda)$ и тогда
        \[
            \phi_X(t) = \Ex e^{itX} = \sum\limits_{k = 0}^{+\infty} e^{ikt} \dfrac{e^{-\lambda}\lambda^k}{k!} = e^{-\lambda}\sum\limits_{k = 0}^{+\infty} \dfrac{(e^{it}\lambda)^k}{k!} = e^{\lambda(e^{it} - 1)}.
        \]
    \end{itemize}
\end{solution}
\begin{problem}
    Вычислить все моменты $\Ex \xi^k, k \in \N$ стандартной нормальной случайной величины и проверить, что её характеристическая функция раскладывается в ряд по моментам:
    \[
        \phi_{\xi}(t) = \sum\limits_{k = 0}^{+\infty} \dfrac{(it)^k}{k!}\Ex(X^k).
    \]
    Вычислить радиус сходимости этого ряда.
\end{problem}
\begin{solution}
    Раньше считались. Радиус сходимости -- $+\infty$.
\end{solution}
\begin{problem}
    Сумма двух независимых целочисленных неотрицательных случайных величин имеет биномиальное распределение.
    Доказать, что каждое слагаемое имеет биномиальное распределение.
\end{problem}
\begin{solution}
    Пусть $X, Y$~---~независимые случайные величины такие, что $X + Y \sim Binom(n, p)$. Тогда
    \[
        \phi_{X + Y}(t) = (1 - p + pe^{it})^n.
    \]
    Но, в то же время, $\phi_{X + Y} = \phi_X \cdot \phi_Y$.
    Поскольку $X$ и $Y$ имеют значения в $\N_0$, то $\phi_X$ и $\phi_Y$ являются степенными рядами от $e^{ikt}, k \in \N$.
    Из единственности разложения в ряд Фурье следует, что $\P(X = k) = 0, k \geq n + 1$.
    Аналогичное верно и про $Y$.
    Тогда получается что $\phi_{X + Y}$ является произведением двух полиномов от $e^{it}$ и следовательно, в силу ОТА, $\phi_X$ и $\phi_Y$ являются $(1 - p + pe^{it})^k$ и $(1 - p + pe^{it})^{n - k}$ что и доказывает требуемое.
\end{solution}
\begin{problem}
    Проверить, являются ли характеристическими функциями каких-нибудь случайных величин функции:
    \begin{enumerate}
        \item $\cos t$;
        \item $|\cos t|$;
        \item $e^{-t^4}$;
        \item $\cos t^2$;
        \item $\dfrac{\sin t}{t}$;
        \item $\re f, \im f$, где $f$~---~характеристическая функция некоторой случайной величины $\xi$.
    \end{enumerate}
\end{problem}
\begin{solution} \leavevmode
    \begin{enumerate}
        \item Заметим, что случайная величина, которая принимает $1$ с вероятностью $1/2$ и $-1$ с вероятностью $1/2$ имеет характеристическую функцию $(e^{it} + e^{-it})/2$ в точности равную косинусу.
        \item Пусть $\psi(t) = |\cos t|$ является характеристической функцией некоторой случайной величины $X$.
        Заметим, что в окрестности нуля функция является просто $\cos x$ (без модуля), а значит является аналитичной и все моменты $X$ существуют и совпадают с моментами случайной величины из предыдущего пункта.
        Но, тогда, ряд Тейлора функции $\psi(t)$ в нуле сходится к ней повсюду поточечно и значит $|\cos t| = \cos t$~---~противоречие. ???
        \item Пусть $\phi(t) = e^{-t^4}$ является характеристической функцией некоторой случайной величины $\xi$.
        Тогда $\Ex \xi^2 = \phi''(0)$. В то же время, $\phi'(t) = -4t^{3}e^{-t^4}$ и $\phi''(0) = -12t^2 e^{-t^4} - 16t^6 e^{-t^4}|_{t = 0} = 0$.
        Значит $\xi^2$ равна нулю почти всюду, а следовательно и $\xi$ равна нулю почти всюду; но у случайной величины, равной нулю почти всюду, характеристическая функция $\phi_\xi = \Ex e^{it\xi} = 1$ по определению, и не равна $e^{-t^4}$.
        \item Пусть $f(t) = \cos t^2$ является характеристической функцией некоторого распределения $\eta$.
        Тогда $\Ex \eta^2 = f''(0)$. Поскольку $f'(t) = -2t \sin t^2$ и $f''(0) = -2\sin t^2 - 4t^2\cos t^2|_{t = 0} = 0$, то по аналогии с предыдущим пунктом мы приходим к противоречию.
        \item Характеристическая функция равномерного распределения на $[0, 1]$ выглядит как $\phi_\xi(t) = \dfrac{e^{it} - 1}{it}$ и $\re \phi_\xi$ в точности равно $\dfrac{\sin t}{t}$.
        \item Пусть $f$~---~характеристическая функция некоторого распределения. Но, тогда $f(0) = 1 \Ra \im f(0) = 0 \neq 1$ и $\im f$ ни при каких условиях не является характеристической функцией распределения.
        Пусть $X$~---~случайная величина, характеристическая функция которой равна $\phi_X$. Тогда, если мы рассмотрим случайную величину $Y$, которая с вероятностью $1/2$ равна $X$ и с такой же вероятностью равна $-X$, то её плотность $p_Y = (p_X + p_{-X})/2$, следовательно в силу линейности $\phi_Y(t) = (\phi_X(t) + \phi_{-X}(t))/2 = (\phi_X(t) + \phi_X(-t))/2 = (\phi_X(t) + \phi_X^+(t))/2 = \re \phi_X$, что и требовалось показать, а следовательно $\re f$~---~всегда характеристическая функция случайной величины.
    \end{enumerate}
\end{solution}
